% !TeX root = ../papers/nc-p3.tex
\section{行列恒等式によらない小次元単純の均衡性と一隅四行列の判定}
\label{sec:generic-corner-small-simples}
\begin{lemma}\label{lem:generic-corner-separation}
標準正則窓の非頂点単純の一般の二倍率 $S,T$ を取る。
ある頂点 $e$ を避ける組成因子だけを $S,T$ の間に入れた表現では、その有限像に
$S$ の行列ブロックから $T$ の行列ブロックへの非零矩形は存在しない。
ここで $S,T$ は各一回で他の因子は頂点単純とする。
\end{lemma}
\begin{proof}
隅 $eBe$ を長さ4を次数1として見ると、有限二次表示を持つ。長さ4の全語をまず生成元に取り、
内部の二次関係を一次式として消して長さ4の経路空間を生成空間にする。
元の二次関係の挿入位置は一つの四文字塊の内部か二つの塊の境界なので、後者は二次関係を与え、
任意の長さでそれらが核を生成する。隅の成長は多項式以下である。
$Se,Te$ は非零単純であり、次数倍率で無限に動く。実際 $Se$ の非零ベクトルは
$S$ 全体を生成し、$e$ を掛け戻せば $Se$ を生成する。さらに $Se$ から $S$ は
$Se\otimes_{eBe}eB$ の唯一の単純商として復元できる。
真部分加群は全て $e$ 成分が零で、その和も真だからである。
定理\ref{thm:growth-generic-extensions}の第一Extの証明はこの隅にも使え、
$\operatorname{Ext}^1_{eBe}(Se,Te)=0$ である。四つの隅に同時に一般の二倍率を選べる。
切断は完全なので $e$ を掛けるとこの二単純しか残らず、三角の向きにかかわらずその拡張は分裂する。
一方、全有限像での非零矩形は左右に行列単位を掛けて全矩形を与え、$e$ と両立する持上げを選んで
$e$ で圧縮しても非零である。この非零な根基成分は分裂表現の隅の像には存在できない。
\end{proof}

\begin{theorem}\label{thm:generic-corner-small-simples}
有限四辺の任意の標準正則窓で、どこか一つの頂点次元が $1,2,3$ の非頂点単純は
全位置でその同じ次元を持つ。二・三次元均衡単純の頂点との第一Extは、
全頂点へ向かう側と全頂点から来る側の少なくとも一方で恒等的に零である。
特に頂点と直交しない二次元単純は次数倍率と同型を除き有限個である。
\end{theorem}
\begin{proof}
次元1は有限点グラフの薄い単純の判定を使う。次元2または3の場合、
定理\ref{thm:opposite-four-matrix-corners}の普遍拡張の不等式を、
そこで「二端だけで超過」とした全強度で適用できる。理由を確かめる。
一般の二倍率間の両向きの第一Extの消滅と $\sum a_i b_i\le E$ は
定理\ref{thm:growth-generic-extensions}による。同証明で作る $U,V$ は、
二端 $S,T$ のほかに $k_i,k_{i+2}$ しか使わず、積の核を貼る三段加群も
これに $k_{i+1}$ を加えるだけである。どれにも $k_{i+3}$ が現れないため、
補題\ref{lem:generic-corner-separation}が以前は恒等式で排除した全ての非零矩形を排除する。
逆向き第一Extの上界は六関係の評価の階数の議論なので同じである。

具体的には前定理の三不等式で $u'=u,w'=w,L_i=L_k=0$ を全位置に採る。
$d_0=2,3$ の各場合、反転まで $2\le d_1\le d_3\le3(d_0-1)$、
$2\le d_2\le3(d_1-1)$ となる。各位置の行・列次元とEuler式を合わせた整数全列挙で
両方とも均衡型だけが残る。\path{computations/verify_generic_corner_small_simples.py}は
前定理の整数計算を強い不等式にしてこの二範囲を再現する。

均衡なら $E=0$、特に頂点ごとに $a_i b_i=0$ である。両方向が異なる位置にあると仮定し、
定理\ref{thm:one-sided-balanced-small-simples}の最短区間での持上げを使う。
その次元上界 $a_i,b_i\le1$ または $\le3$ は有限辺だけから出る。そこでの経路汎関数 $w$ と
次数付き中間加群 $N_w$ は一対の一般倍率について選べばよい。
$N_w$ は次数倍率で同型を保つため、一旦二端で必要な拡張を得れば、二つの倍率を
\emph{独立に}動かしても同じ $N_w$ で両端の非分裂制限を保てる。
末尾からの制限と頭への押し出しもその次数付き同型で保たれる。
一般の二端間の第二Extは零で、三段像には頭から末尾を通る非零積がある。
中間の全頂点を経る長さ3の場合も含めて補題\ref{lem:growth-fixed-middle}に反する。
従って全体で高々片方向となり、最後の有限性は補題\ref{lem:binary-simple-trace-escape}から従う。
\end{proof}

\begin{corollary}\label{cor:single-four-matrix-all-regular}
四辺有限な標準有限Frobenius候補で、標準Koszul左右AS正則かつ\emph{一隅}が四行四列行列の
全恒等式を満たすことは、定理\ref{thm:four-matrix-all-regular}の八次中心元による32次までの
判定と一般隅の指数検査に同値である。全ての該当窓は中心有限な合成整域で指数は2か4である。
\end{corollary}
\begin{proof}
指定隅の全単純の次元は4以下で折り畳みはPIである。指定隅の次元4なら中心多項式を延ばして
頂点に直交し均衡する。次元3以下の均衡性と片側消滅は直前の定理による。
片側だけが残る三次元単純に薄い $D$ を使うと、補題\ref{lem:periodic-tail-vertex-bridges}は
指定隅で $1+3+1>4$ に反する。直交する非薄単純が一つでもあるときは
同じ補題で二次元単純の残る方向も消える。従ってこの場合、全非頂点単純が直交して
定理\ref{thm:pi-vertex-extension-criterion}で中心有限となる。直交する非薄単純がなければ、
前定理の有限性から全単純が有限個の次数倍率軌道となる。
定理\ref{thm:four-matrix-all-regular}の証明のPI根基の冪零性とAS消滅による矛盾がそのまま使える。
中心有限の場合の合成整性と全隅の指数の一致により、指定隅の次数上界は全隅に及ぶ。
残る同値と指数は同定理による。
\end{proof}
